\documentclass[11pt,a4paper]{article}
\usepackage[utf8]{inputenc}
\usepackage{amsmath,amssymb,amsthm}
\usepackage{graphicx}
\usepackage{booktabs}
\usepackage{hyperref}
\usepackage{natbib}
\usepackage{geometry}
\usepackage{xcolor}
\usepackage{float}
\usepackage{enumitem}

\geometry{margin=1in}

\newtheorem{definition}{Definition}
\newtheorem{theorem}{Theorem}
\newtheorem{proposition}{Proposition}
\newtheorem{remark}{Remark}

\newcommand{\E}{\mathbb{E}}
\newcommand{\Prob}{\mathbb{P}}
\newcommand{\LCDM}{$\Lambda$CDM}

\title{An E-Value Perspective on DESI DR2 Evidence for Dynamical Dark Energy}

\author{
Jinyoung Kim$^{1}$\thanks{Correspondence: \texttt{jinyoungkim927@gmail.com}} \and
Andrius Tamosiunas$^{2}$ \and
David F. Mota$^{2}$\\[0.5em]
\small $^{1}$Department of Computer Science, Stanford University\\
\small $^{2}$Institute of Theoretical Astrophysics, University of Oslo
}

\date{April 2026}

\begin{document}

\maketitle

\begin{abstract}
The DESI collaboration recently reported 3--4$\sigma$ frequentist evidence for dynamical dark energy ($w_0w_a$CDM) over the cosmological constant (\LCDM). This note applies e-value methodology---a recently developed framework from the statistics literature that provides a notion of evidence with a sample-size-free, $\alpha$-free frequentist interpretation---to the publicly available DESI DR2 BAO data. Our most informative finding concerns cross-dataset consistency: DES-Y5's best-fit $w_0w_a$CDM parameters predict DESI data \textit{worse} than \LCDM{} ($E = 0.19$), while Pantheon+'s parameters predict it well ($E = 2049$)---a ${\sim}10{,}000\times$ asymmetry pointing to tension between supernova catalogs rather than a coherent dark energy signal. Regarding the DESI data alone, two independent valid e-value methods converge on moderate evidence: the uniform mixture e-value yields $E \approx 15$ ($\ln E \approx 2.7$) and the leave-one-out (LOO) average yields $E \approx 10$ ($\ln E \approx 2.3$). A data-split e-value gives $E = 1.4$, though power calibration shows this test has limited statistical power for $w_a$ (median $E \approx 2.7$ even when $w_0w_a$CDM is the true model). These results, together with recent Bayesian analyses by \citet{OngBayesian2025} and tension diagnostics by \citet{WangMota2025}, suggest that the interpretation of DESI's signal depends substantially on methodology and dataset combination, and that resolving inter-dataset tensions is a prerequisite for strong conclusions about dark energy dynamics.
\end{abstract}

\vspace{1em}
\noindent\textbf{Note on Authorship:} An initial draft of this manuscript was prepared by JK with assistance from a large language model (Claude). The research direction, methodology, and interpretation were developed by the authors, who have reviewed, edited, and approved all content. The authors welcome feedback on both the statistical methodology and its application to cosmological data. (To be revised in line with the AI-use policy of the chosen target journal.)

\tableofcontents
\newpage

%=====================================================================
\section{Introduction}
\label{sec:intro}
%=====================================================================

The standard cosmological model, $\Lambda$ cold dark matter (\LCDM), describes a spatially flat universe whose late-time accelerated expansion is driven by a cosmological constant $\Lambda$ with equation-of-state parameter $w = -1$. Over the past two decades \LCDM{} has provided a remarkably successful fit to cosmic microwave background (CMB) anisotropies, type Ia supernovae, weak lensing, galaxy clustering, and baryon acoustic oscillations (BAO) with only six free parameters. The same data sets have, however, begun to expose tensions that the model is increasingly hard pressed to absorb---most prominently the $\sim 5\sigma$ discrepancy between local distance-ladder and CMB-inferred determinations of the Hubble constant $H_0$, and the persistent $S_8$ tension between primary CMB and weak-lensing measurements of the late-time clustering amplitude. Whether these tensions reflect systematics, statistical fluctuations, or new physics is unresolved, and a large literature has explored modified-gravity, interacting-dark-sector, and early-dark-energy extensions of the model.

Against this backdrop, the Dark Energy Spectroscopic Instrument (DESI) collaboration's recent BAO measurements---the most precise of their kind---have attracted considerable attention. Combined with CMB and supernova data, the DR2 release reported a $\sim 3$--$4\sigma$ frequentist preference for a time-evolving dark energy equation of state ($w_0w_a$CDM) over the cosmological constant \citep{DESI2025}. If taken at face value, this would be the first statistically significant evidence for dynamical dark energy and a substantive departure from \LCDM. The interpretation of this signal is, however, contested: several recent analyses using complementary statistical approaches have argued that the significance is more fragile than the headline number suggests, and may be driven in part by tensions \emph{between} the datasets combined to produce it rather than by a coherent physical signal. The remainder of this section reviews these analyses and motivates the e-value perspective developed below.

\noindent\textbf{[Note for AT: placeholder draft of the broader astro framing; please rewrite as you see fit per your Overleaf comment.]}

\subsection{Context and Motivation}

The DESI DR2 collaboration reported a $\chi^2$ improvement of approximately 12 when fitting the $w_0w_a$CDM model (with two additional parameters $w_0$ and $w_a$) compared to the standard \LCDM{} model using their BAO measurements combined with CMB and supernova data \citep{DESI2025}. Interpreted through standard frequentist methodology, this corresponds to 2.5--4$\sigma$ significance favoring dynamical dark energy.\footnote{The precise significance depends on the degrees of freedom used for the $\sigma$-conversion. For $k=2$ extra parameters, $\Delta\chi^2 = 11.9$ should be evaluated against the $\chi^2$ distribution with 2 d.o.f.: $p = 1 - F_{\chi^2}(11.9;\, 2) = 0.0026$, corresponding to $\approx 2.8\sigma$ (two-sided). The sometimes-quoted ``$3.5\sigma$'' uses $\sqrt{\Delta\chi^2}$, which is valid only for $k=1$. DESI's reported range of 3--4$\sigma$ reflects different dataset combinations and methodological choices.}

Two recent independent analyses have questioned this interpretation:
\begin{itemize}
    \item \citet{OngBayesian2025} performed Bayesian model comparison using nested sampling, finding that for DESI+CMB data, the Bayes factor modestly favors \LCDM{} ($\ln \mathcal{B} = -0.57 \pm 0.26$), though this result depends on the choice of prior over $(w_0, w_a)$. They identified a significant tension (2.95$\sigma$) between DESI and DES-Y5 supernovae within \LCDM{}, which $w_0w_a$CDM resolves.

    \item \citet{WangMota2025} documented parameter tensions between CMB, DESI, and various supernova catalogs, concluding that combined constraints may be problematic due to dataset inconsistencies.
\end{itemize}

This note contributes a third perspective using \textbf{e-values}, a framework from the recent statistics literature \citep{Vovk2021,Shafer2021,Ramdas2023} that addresses concerns about post-hoc model selection and overfitting. We apply this methodology to examine whether the apparent preference for $w_0w_a$CDM generalizes out-of-sample.

\subsection{Scope and Limitations}

We emphasize several limitations of this analysis:
\begin{enumerate}
    \item We analyze the publicly released BAO summary statistics, not the underlying galaxy catalogs. Our conclusions are conditional on the accuracy of the published data products.

    \item Our cosmological calculations use standard distance formulas with Planck 2018 fiducial parameters. More sophisticated treatments (e.g., full MCMC with CAMB/CLASS) may yield quantitatively different results.

    \item E-values provide one valid perspective on hypothesis testing but are not universally superior to other approaches. Different methods answer different questions.

    \item The lead author's primary background is in statistics and computer science rather than cosmology, with the cosmological framing and astrophysical interpretation contributed by the Oslo co-authors. This work should therefore be read as an application of statistical methodology to published data, undertaken jointly with cosmologists, rather than as a comprehensive new cosmological analysis.
\end{enumerate}

\subsection{Summary of Approach}

Our analysis proceeds as follows:
\begin{enumerate}
    \item We briefly define e-values and their key property (Section~\ref{sec:evalues}), with formal statements in Appendix~\ref{app:proofs}.

    \item We describe the DESI data and our methodology (Section~\ref{sec:methods}).

    \item We present results, leading with the cross-dataset tension analysis and followed by within-dataset tests (Section~\ref{sec:results}).

    \item We discuss how these findings relate to existing literature (Section~\ref{sec:discussion}).
\end{enumerate}

%=====================================================================
\section{E-Values: A Brief Overview}
\label{sec:evalues}
%=====================================================================

\subsection{Definition and Interpretation}

Let $H_0$ denote the null hypothesis (here, \LCDM, in which $w = -1$ exactly), let $H_1$ denote the alternative hypothesis ($w_0w_a$CDM, with $w(a) = w_0 + w_a(1-a)$), and let $\E[\,\cdot\, \mid H_0]$ denote expectation under the data-generating distribution implied by $H_0$. An \textbf{e-value} is a non-negative random variable $E$ satisfying $\E[E \mid H_0] \leq 1$. The concept has deep roots in the work of Ville (1939) on martingale theory, but has been rediscovered and systematized in recent years by \citet{Vovk2021}, \citet{Shafer2021}, and \citet{Ramdas2023} as a unifying framework for hypothesis testing. Following \citet{GrunwaldBeyondNeymanPearson}, we emphasize that e-values provide a notion of evidence with a frequentist meaning that does not require a significance level $\alpha$ to be fixed in advance and that does not depend on sample size: a value of $E$ can be reported and interpreted on its own terms, in addition to being convertible into a Type~I-controlled test via Ville's inequality. For cosmologists familiar with Bayesian and frequentist methods, e-values can be understood as follows: like a Bayes factor, an e-value quantifies relative evidence between hypotheses; like a $p$-value, it provides Type~I error control (via Ville's inequality rather than tail probabilities). Unlike Bayes factors, e-values do not require a full prior distribution---a mixture over a discrete grid suffices. Unlike $p$-values, e-values can be multiplied across independent datasets (enabling sequential analysis) and remain valid under optional stopping.

The key properties are:
\begin{itemize}
    \item \textbf{Type I error control:} Rejecting $H_0$ when $E \geq 1/\alpha$ yields a test with significance level $\alpha$ (see Theorem~\ref{thm:ville} in the Appendix).

    \item \textbf{Combination:} Independent e-values can be multiplied to accumulate evidence.

    \item \textbf{Interpretation:} An e-value of $E$ can be thought of as ``the data are $E$ times more consistent with the alternative than expected under the null.''
\end{itemize}

\subsection{Likelihood Ratios and the Overfitting Problem}

For simple hypotheses (where both $H_0$ and $H_1$ are fully specified before seeing the data), the likelihood ratio $E = L(\text{data} \mid H_1) / L(\text{data} \mid H_0)$ is a valid e-value. However, when $H_1$ involves parameters fitted to the same data used for testing---i.e., plugging in the maximum likelihood estimator (MLE)---the result is \textbf{not an e-value at all}. Specifically, under $H_0$ the $\Delta\chi^2$ for $k$ extra parameters follows a $\chi^2(k)$ distribution, so the maximized likelihood ratio has expectation
\begin{equation}
\E_{H_0}\!\left[\exp\!\left(\tfrac{\chi^2(k)}{2}\right)\right] = \int_0^\infty e^{t/2}\,\frac{t^{k/2-1}e^{-t/2}}{2^{k/2}\Gamma(k/2)}\,dt = \infty \quad \text{for } k \geq 2,
\end{equation}
where $\Gamma(\cdot)$ is the Gamma function. For the $w_0w_a$CDM case ($k=2$ extra parameters), one can verify directly:
\begin{equation}
\E_{H_0}\!\left[e^{\chi^2(2)/2}\right] = \int_0^\infty e^{t/2}\cdot\tfrac{1}{2}e^{-t/2}\,dt = \int_0^\infty \tfrac{1}{2}\,dt = \infty.
\end{equation}
Since $\E[E \mid H_0] = \infty$, the defining property $\E[E \mid H_0] \leq 1$ is violated: \textbf{the maximized likelihood ratio is not a valid e-value}. It should be understood as a descriptive statistic, not as evidence with a frequentist guarantee. We include it in Section~\ref{sec:samedata} only for comparison to show how drastically the evidence changes under valid methods.

\subsection{Data-Split E-Values}

To construct a valid e-value when the alternative involves fitted parameters, we use \textbf{data splitting}:
\begin{enumerate}
    \item Split data into training set $D_{\text{train}}$ and test set $D_{\text{test}}$
    \item Fit parameters $\hat{\theta}$ using only $D_{\text{train}}$
    \item Compute $E = L(D_{\text{test}} \mid \hat{\theta}) / L(D_{\text{test}} \mid H_0)$
\end{enumerate}

Because $\hat{\theta}$ is determined before observing $D_{\text{test}}$, this yields a valid e-value (Proposition~\ref{prop:split} in the Appendix). The trade-off is reduced statistical power due to using only part of the data for testing.

\subsection{Uniform Mixture E-Values}

An alternative approach averages the likelihood ratio over a grid of alternative parameter values:
\begin{equation}
E_{\text{mix}} = \frac{1}{|\mathcal{G}|} \sum_{(w_0, w_a) \in \mathcal{G}} \frac{L(\text{data} \mid w_0, w_a)}{L(\text{data} \mid H_0)}
\end{equation}
where $\mathcal{G}$ is a uniform grid over the $(w_0, w_a)$ space. Because each term in the sum is a likelihood ratio for a pre-specified (simple) alternative, and the average of e-values is an e-value by linearity of expectation, $E_{\text{mix}}$ is a valid e-value (Proposition~\ref{prop:mixture} in the Appendix). This method uses all the data but pays an Occam penalty: probability mass ``wasted'' on parameter values far from the truth reduces the e-value relative to the maximized likelihood ratio.

\subsection{Leave-One-Out (LOO) Average E-Values}

We also construct e-values by leave-one-out cross-validation across the 7 redshift bins. For each bin $k$:
\begin{enumerate}
    \item Fit $(w_0, w_a)$ on the remaining 6 bins (training set)
    \item Compute the likelihood ratio $E_k$ on the held-out bin $k$
\end{enumerate}
Each individual $E_k$ is a valid e-value (conditional on the training data, the held-out bin tests a pre-specified alternative). The \textbf{average} $\bar{E} = \frac{1}{K}\sum_{k=1}^K E_k$ is a valid e-value by linearity of expectation (Proposition~\ref{prop:loo} in the Appendix): $\E[\bar{E} \mid H_0] = \frac{1}{K}\sum_k \E[E_k \mid H_0] \leq 1$.

\paragraph{The LOO product is not a valid e-value.} The product $\prod_k E_k$ does not in general satisfy $\E[\prod_k E_k \mid H_0] \leq 1$. Multiplication preserves the e-value property only when the factors are independent, and the leave-one-out training sets overlap heavily by construction (any two folds share $K-2$ of $K$ bins), so the fitted parameters $\hat{\theta}_{-k}$ are strongly correlated across $k$. The expectation of the LOO product under $H_0$ has no closed form and may exceed unity; we therefore report only the LOO average $\bar{E}$, which retains validity by linearity of expectation.

\subsection{Power-Optimized Point-Alternative E-Values}
\label{sec:ump_theory}

The methods above---data splitting, uniform mixtures, and LOO averaging---all sacrifice power in different ways: data splitting discards test data, mixtures dilute across a grid, and LOO averaging divides by $K$. A natural question is whether one can construct an e-value that uses all the data and avoids the averaging penalty while remaining valid. The general framework of \citet{GrunwaldSafeTesting2024} provides one such construction, with two important caveats that we make explicit below.

For a \textit{simple} null $H_0$ (such as \LCDM, where $w = -1$ exactly), the likelihood ratio $E = L(\text{data} \mid \theta_1) / L(\text{data} \mid H_0)$ for any pre-specified point alternative $\theta_1$ is a valid e-value by Proposition~\ref{prop:lr}. One may therefore choose $\theta_1$ to maximize statistical power without compromising validity, \emph{provided} $\theta_1$ is determined before any contact with the test data and depends only on auxiliary information (the covariance structure, simulations from a class of plausible alternatives, etc.). We refer to the resulting object as a \textbf{power-optimized point-alternative e-value}; it is sometimes called a UMP e-value in the literature, though strictly speaking the UMP terminology applies only in restricted settings (e.g.\ one-sided tests in single-parameter exponential families) and for the present two-parameter composite alternative the construction is better thought of as a power-targeted simple-vs-simple test.

Concretely, fix a significance level $\alpha$ (e.g., $\alpha = 0.01$). For each candidate $\theta_1 = (w_0, w_a)$, the rejection region is $\{x : L(x \mid \theta_1)/L(x \mid H_0) \geq 1/\alpha\}$. We select $\theta_1^*$ to maximize the rejection probability under a pre-specified ensemble of ``representative'' alternatives via Monte Carlo simulation (a maximin criterion across the ensemble). The resulting $\theta_1^*$ is then fixed, and the likelihood ratio on the real data is computed. Because $\theta_1^*$ depends only on the covariance structure and the simulation design---not on the observed data---the result is a valid e-value regardless of $\alpha$.

\paragraph{Two important caveats.} First, the construction is genuinely simple-vs-simple: the resulting e-value tests \LCDM{} against the single point $\theta_1^*$, not against the full composite $w_0w_a$CDM model. A high value of this e-value tells us that the data look much more consistent with this one specific $(w_0, w_a)$ than with \LCDM, but it does not by itself answer the broader question of how strongly the data favor the composite alternative as a class. Second, the choice of $\theta_1^*$ depends on the ensemble of ``true'' alternatives used in the power simulation. Different ensembles (or differently weighted maximin criteria) generally select different $\theta_1^*$, and hence yield different e-values on the real data. Unlike the uniform mixture, which makes its averaging explicit, the power-optimized construction conceals an analogous design choice inside the simulation step. We return to both points when interpreting our numerical results in Section~\ref{sec:ump}.

%=====================================================================
\section{Data and Methods}
\label{sec:methods}
%=====================================================================

\subsection{DESI DR2 BAO Data}

We use the official DESI DR2 BAO measurements from the CobayaSampler repository,\footnote{\url{https://github.com/CobayaSampler/bao_data}} comprising 13 measurements across 7 redshift bins ($z_{\text{eff}} = 0.295$ to $z_{\text{eff}} = 2.33$). The seven tracer samples are the Bright Galaxy Sample (BGS, $z_{\text{eff}} = 0.295$, $D_V/r_d$), three Luminous Red Galaxy bins (LRG1, $z_{\text{eff}} = 0.510$; LRG2, $z_{\text{eff}} = 0.706$; combined LRG3 + Emission Line Galaxy bin ELG1, $z_{\text{eff}} = 0.934$), a second Emission Line Galaxy bin (ELG2, $z_{\text{eff}} = 1.317$), Quasi-Stellar Objects (QSO, $z_{\text{eff}} = 1.491$), and the Lyman-$\alpha$ forest (Ly$\alpha$, $z_{\text{eff}} = 2.330$), where the last six each provide both the transverse comoving distance $D_M/r_d$ and the Hubble distance $D_H/r_d$. The full $13\times 13$ covariance matrix is block-diagonal by redshift bin: the off-diagonal blocks between different bins are zero, while within each bin the $D_M/r_d$ and $D_H/r_d$ measurements are anti-correlated with correlation coefficients ranging from $-0.4$ to $-0.5$ (see Fig.~\ref{fig:correlation}). All values are taken directly from arXiv:2503.14738, Table~IV \citep{DESI2025}.

Figure~\ref{fig:bao_data} shows the 13 BAO measurements with their $1\sigma$ error bars, overlaid with theoretical predictions for \LCDM{} (solid) and $w_0w_a$CDM with $(w_0, w_a) = (-0.75, -1.05)$ (dashed). Figure~\ref{fig:correlation} displays the correlation matrix, confirming the block-diagonal structure: bins are uncorrelated, while $D_M/r_d$ and $D_H/r_d$ within each bin are anti-correlated.

\begin{figure}[H]
    \centering
    \includegraphics[width=\textwidth]{figure_bao_data.pdf}
    \caption{DESI DR2 BAO measurements with $1\sigma$ error bars. Left: volume-averaged distance $D_V/r_d$ (BGS only). Center: transverse comoving distance $D_M/r_d$. Right: Hubble distance $D_H/r_d$. Solid black curves show the \LCDM{} prediction $(w_0 = -1,\, w_a = 0)$; dashed red curves show the $w_0w_a$CDM best fit $(w_0 = -0.75,\, w_a = -1.05)$. Both curves use the same Planck 2018 background parameters $h = 0.6766$, $\Omega_m = 0.3111$, $r_d = 147.05\,$Mpc; only $(w_0, w_a)$ differ between the two models. Tracer labels indicate the galaxy sample at each effective redshift.}
    \label{fig:bao_data}
\end{figure}

\begin{figure}[H]
    \centering
    \includegraphics[width=0.75\textwidth]{figure_correlation.pdf}
    \caption{Correlation matrix of the 13 DESI DR2 BAO measurements. The matrix is block-diagonal by redshift bin: off-diagonal blocks are exactly zero (bins are independent), while within each bin, $D_M/r_d$ and $D_H/r_d$ show anti-correlations of $-0.4$ to $-0.5$. Block outlines highlight the redshift-bin structure.}
    \label{fig:correlation}
\end{figure}

\subsection{Models Compared}
\label{sec:models}

\textbf{Null hypothesis ($H_0$):} \LCDM{} with $w = -1$ (cosmological constant).

\textbf{Alternative hypothesis ($H_1$):} $w_0w_a$CDM with the Chevallier--Polarski--Linder (CPL) equation of state $w(a) = w_0 + w_a(1-a)$ \citep{Chevallier2001,Linder2003}, where $a = 1/(1+z)$ is the cosmological scale factor and $w_0$, $w_a$ are free parameters.

\paragraph{Fixed cosmological parameters.} In both models, we fix the background cosmological parameters to their Planck 2018 best-fit values: $h = 0.6766$, $\Omega_m = 0.3111$, and $r_d = 147.05\,$Mpc. Only $w_0$ and $w_a$ are varied when fitting the alternative hypothesis.

This is standard practice for BAO-only analyses and is appropriate here because BAO summary statistics ($D_M/r_d$, $D_H/r_d$, $D_V/r_d$) depend on $\{h, \Omega_m, r_d\}$ and $\{w_0, w_a\}$, but the former are tightly constrained by CMB data independently of dark energy dynamics. The remaining \LCDM{} parameters ($\Omega_b h^2$, $n_s$, $A_s$, $\tau$) enter BAO predictions only through $r_d$, which we hold fixed. Consequently, the effective number of additional parameters in $w_0w_a$CDM relative to \LCDM{} is $k = 2$ (i.e., $w_0$ and $w_a$), not the full parameter count of each model.

In principle, one could marginalize over $\{h, \Omega_m, r_d\}$ within the e-value framework---for instance, by extending the uniform mixture grid to include these parameters, or by fitting them jointly with $(w_0, w_a)$ in the data-split and LOO procedures. However, BAO data alone cannot simultaneously constrain $\Omega_m$, $h$, and $r_d$ without strong degeneracies, so informative priors (effectively, CMB constraints) would be required. We leave such an extended analysis for future work and note this as a limitation in Section~\ref{sec:limitations}.

\paragraph{Source of best-fit $(w_0, w_a)$ values.} For the cross-dataset analysis (Section~\ref{sec:cross}), we use published best-fit $(w_0, w_a)$ values from supernova-only analyses. For the data-split and LOO analyses, $(w_0, w_a)$ are fitted directly to the DESI training data via $\chi^2$ minimization. For the power calibration, we use DESI's BAO-only best-fit $(w_0 = -0.75, w_a = -1.05)$ as the true model. These values are approximate, drawn from various tables in \citet{DESI2025} and the associated supernova literature; the precise values depend on which dataset combination and methodology are used.

\subsection{Cross-Dataset Validation}

We examine whether $(w_0, w_a)$ values fitted on one experiment predict another experiment better than \LCDM. We use published best-fit values from the three supernova compilations \citep{DESI2025}:
\begin{itemize}
    \item Pantheon+: $w_0 \approx -0.90$, $w_a \approx -0.20$
    \item Union3: $w_0 \approx -0.78$, $w_a \approx -0.80$
    \item DES-Y5: $w_0 \approx -0.65$, $w_a \approx -1.20$
\end{itemize}
We list these values as approximate (``$\approx$'') for two reasons. First, they are read off published summary tables and figures rather than reproduced from the underlying chains: each compilation reports slightly different best-fit values depending on whether the result quoted is the maximum a posteriori, the posterior mean, or the marginalized 1D mode, and depending on which DESI/CMB combination is included alongside the supernovae. Second, the values themselves shift at the $\sim 0.01$--$0.02$ level between different DESI papers and supernova releases as systematic treatments are updated. Because the cross-dataset e-value depends only weakly on $(w_0, w_a)$ at this level (and the headline asymmetry between Pantheon+ and DES-Y5 is many orders of magnitude), the approximation is harmless for the qualitative conclusions of Section~\ref{sec:cross}; a fully posterior-marginalized version is discussed as a refinement in Section~\ref{sec:limitations}.

\subsection{Data-Split Procedure}

We split the 13 DESI measurements by redshift at $z = 1$:
\begin{itemize}
    \item \textbf{Training:} 7 measurements at $z < 1$ (BGS, LRG1, LRG2, LRG3+ELG1)
    \item \textbf{Test:} 6 measurements at $z \geq 1$ (ELG2, QSO, Ly$\alpha$)
\end{itemize}

The $z = 1$ threshold is chosen because it approximately bisects the data (7 vs.\ 6 measurements) and provides a natural physical boundary: the low-$z$ training set constrains the equation of state near the present epoch, while the high-$z$ test set probes earlier times where dark energy effects differ most between models. We also test alternative splits (alternating bins, random 50/50) in Table~\ref{tab:splits} to assess sensitivity.

We fit $(w_0, w_a)$ by minimizing $\chi^2$ on the training set, then evaluate the likelihood ratio on the test set.

\subsection{Leave-One-Out Procedure}

For the LOO analysis, we iterate over the 7 redshift bins (BGS, LRG1, LRG2, LRG3+ELG1, ELG2, QSO, Ly$\alpha$). In each fold $k$, we fit $(w_0, w_a)$ on the 6 remaining bins and compute the likelihood ratio on the held-out bin. We report both the individual per-bin e-values and the average $\bar{E} = \frac{1}{7}\sum_{k=1}^7 E_k$.

\subsection{Power Calibration}

To assess the statistical power of the data-split test, we run 500 Monte Carlo simulations. In each simulation, we compute BAO distance predictions under $w_0w_a$CDM with parameters $(w_0 = -0.75, w_a = -1.05)$ and all other cosmological parameters fixed to their fiducial values (same as in the main analysis; see Section~\ref{sec:models}). We then add Gaussian noise drawn from the published DESI covariance matrix to generate a synthetic dataset, apply the same redshift-based data split ($z = 1$ threshold), and compute the resulting e-value. The distribution of $E_{\text{split}}$ under this alternative gives us the expected power of the test.

\subsection{Assumptions}

Our analysis assumes:
\begin{enumerate}
    \item The published DESI covariance matrix accurately captures measurement uncertainties.
    \item Different redshift bins are approximately independent (the covariance matrix is block-diagonal).
    \item The cosmological distance calculations are sufficiently accurate for this comparison.
    \item The CPL parameterization $w(a) = w_0 + w_a(1-a)$ adequately captures potential dark energy dynamics.
\end{enumerate}

%=====================================================================
\section{Results}
\label{sec:results}
%=====================================================================

\subsection{Cross-Dataset E-Values: The Supernova Catalog Asymmetry}
\label{sec:cross}

We begin with what we consider the most informative result: cross-dataset e-values, which test whether $w_0w_a$CDM parameters fitted on one experiment predict another experiment better than \LCDM. Unlike within-dataset tests such as data splitting, this approach is not vulnerable to the objection that low e-values merely reflect insufficient statistical power. The logic is as follows: if two independent experiments both detect the same underlying physics, the best-fit parameters from one \textit{must} predict the other well, because they are measuring the same equation of state $w(a) = w_0 + w_a(1-a)$. A failure of cross-prediction ($E < 1$) therefore constitutes positive evidence of tension between datasets, not merely an absence of evidence for new physics.

Table~\ref{tab:cross} presents the cross-dataset e-values.

\begin{table}[H]
\centering
\caption{Cross-dataset e-values: parameters from one dataset predicting another}
\label{tab:cross}
\begin{tabular}{@{}llcc@{}}
\toprule
Training Dataset & Test Dataset & $(w_0, w_a)$ & E-value \\
\midrule
DESI (fitted) & Pantheon+ & $(-0.86, -0.43)$ & 1.5 \\
DESI (fitted) & DES-Y5 & $(-0.86, -0.43)$ & 86 \\
Pantheon+ & DESI & $(-0.90, -0.20)$ & 2049 \\
DES-Y5 & DESI & $(-0.65, -1.20)$ & \textbf{0.19} \\
\bottomrule
\end{tabular}
\end{table}

\textbf{The central finding:} DES-Y5's best-fit $w_0w_a$CDM parameters ($w_0 = -0.65$, $w_a = -1.20$) predict the DESI BAO data \textit{worse} than \LCDM{} ($E = 0.19 < 1$). Meanwhile, Pantheon+'s best-fit parameters ($w_0 = -0.90$, $w_a = -0.20$) predict DESI data well ($E = 2049$). This represents a factor of ${\sim}10{,}000$ asymmetry between the two leading supernova catalogs in their ability to predict DESI under the $w_0w_a$CDM framework.

The interpretation is straightforward. If $w_0w_a$CDM represents real physics, all experiments are measuring the same underlying equation of state, and the best-fit $(w_0, w_a)$ from one experiment should predict another at least as well as \LCDM, error-bar differences notwithstanding. That DES-Y5's parameters make DESI data \textit{less} probable than \LCDM{} ($E < 1$) means DES-Y5 and DESI prefer incompatible regions of $(w_0, w_a)$ space---a direct detection of inconsistency, not a power deficit.

This ${\sim}10{,}000\times$ asymmetry aligns with two independent analyses: \citet{OngBayesian2025} identified a 2.95$\sigma$ tension between DESI and DES-Y5 within \LCDM{} that $w_0w_a$CDM resolves by absorbing into its additional parameters, and \citet{WangMota2025} documented parameter inconsistencies across CMB, DESI, and supernova catalogs more broadly. The dataset tension itself was identified by those authors; our contribution is the cross-prediction e-value framework, which produces a single number quantifying how well one experiment's preferred dark energy model predicts another's data. $E = 0.19$ measures not merely a parameter discrepancy but an out-of-sample predictive failure.

\begin{figure}[H]
\centering
\includegraphics[width=0.75\textwidth]{figure6_cross_dataset.pdf}
\caption{Cross-dataset e-value matrix. Green indicates the training parameters predict the test data better than \LCDM; red indicates worse. The ${\sim}10{,}000\times$ asymmetry between Pantheon+ ($E = 2049$) and DES-Y5 ($E = 0.19$) when predicting DESI reveals fundamental tension between supernova catalogs within the $w_0w_a$CDM framework.}
\label{fig:cross}
\end{figure}

\subsection{Data-Split E-Value}
\label{sec:datasplit}

We next examine whether the signal generalizes within the DESI data itself. Fitting on the training set ($z < 1$) yields $\hat{w}_0 = -0.78$, $\hat{w}_a = -0.52$. Evaluating on the held-out test set ($z \geq 1$):
\begin{equation}
\boxed{E_{\text{split}} = 1.4}
\end{equation}

\textbf{Interpretation:} The fitted $w_0w_a$CDM model predicts the high-redshift data only 1.4 times better than \LCDM. However, this result must be interpreted with an important caveat about statistical power.

\paragraph{Reduced power for $w_a$.} The data-split test has limited sensitivity to $w_a$ specifically. In the CPL parameterization, $w_a$ contributes $(1-a)\,w_a$ to the equation of state; for the training set at $z < 1$ (i.e.\ $a > 0.5$), the factor $(1-a)$ ranges from $\sim 0.23$ to $0.48$, so the training data have only fractional leverage on $w_a$. Since much of DESI's claimed signal lives in the $w_a$ degree of freedom, a test that poorly constrains $w_a$ from the training set will yield a weak e-value regardless of whether the signal is real. Hence $E = 1.4$ is \textit{inconclusive} rather than negative---consistent both with \LCDM{} and with a genuine $w_0w_a$CDM signal that the test lacks the power to detect---which is why we lead with the cross-dataset result of Section~\ref{sec:cross}.

\paragraph{Power calibration.} Monte Carlo simulation (500 trials) under the DESI best-fit $w_0w_a$CDM model shows that the median data-split e-value is only $E \approx 2.7$ even when $w_0w_a$CDM is the true model. The observed $E = 1.4$ falls within the interquartile range of the distribution under $H_1$. This confirms that $E = 1.4$ cannot distinguish between ``the signal is absent'' and ``the test lacks power to detect it.''

\begin{figure}[H]
\centering
\includegraphics[width=0.85\textwidth]{figure1_evalue_comparison.pdf}
\caption{Comparison across methods. The same-data maximized likelihood ratio (red) is not a valid e-value ($\E[E|H_0]=\infty$). Data-split validation (green) is a valid e-value and yields $E = 1.4$.}
\label{fig:evalue}
\end{figure}

\subsection{Leave-One-Out Average E-Value}
\label{sec:loo}

Table~\ref{tab:loo} presents the per-bin LOO e-values.

\begin{table}[H]
\centering
\caption{Leave-one-out e-values by redshift bin. The redshift $z$ listed is the effective redshift $z_{\text{eff}}$ of each bin (not the bin edges).}
\label{tab:loo}
\begin{tabular}{@{}lccc@{}}
\toprule
Held-out bin & $z_{\text{eff}}$ & Fitted $(w_0, w_a)$ & $E_k$ \\
\midrule
BGS & 0.295 & $(-0.78, -0.82)$ & 1.89 \\
LRG1 & 0.510 & $(-0.80, -0.75)$ & 0.71 \\
LRG2 & 0.706 & $(-0.69, -1.18)$ & 55.98 \\
LRG3+ELG1 & 0.934 & $(-0.93, -0.18)$ & 8.73 \\
ELG2 & 1.317 & $(-0.80, -0.71)$ & 2.14 \\
QSO & 1.491 & $(-0.83, -0.66)$ & 0.75 \\
Ly$\alpha$ & 2.330 & $(-0.85, -0.54)$ & 1.00 \\
\midrule
\multicolumn{3}{@{}l}{\textbf{LOO average} $\bar{E} = \frac{1}{7}\sum E_k$} & \textbf{10.2} \\
\multicolumn{3}{@{}l}{LOO product $\prod E_k$ (\textit{not valid})} & \textit{1062} \\
\bottomrule
\end{tabular}
\end{table}

The LOO average $\bar{E} \approx 10.2$ ($\ln \bar{E} \approx 2.3$) constitutes moderate evidence for $w_0w_a$CDM. The evidence is concentrated in LRG2 ($z = 0.706$, $E_k = 55.98$) and LRG3+ELG1 ($z = 0.934$, $E_k = 8.73$), which are the bins where DESI departs most from \LCDM{} predictions.

\paragraph{The LOO product is reported only for reference.} The product $\prod_k E_k = 1062$ does not constitute calibrated evidence: as discussed in Section~\ref{sec:evalues}, the overlapping training sets violate the independence condition required for $\E\!\big[\prod_k E_k\big] = \prod_k \E[E_k]$ to hold. We list the value purely so that readers can recover it from the per-bin entries; only the LOO average $\bar{E}$ should be interpreted as evidence.

\subsection{Uniform Mixture E-Values and Sensitivity}
\label{sec:mixture}

\begin{table}[H]
\centering
\caption{E-values under different data splits}
\label{tab:splits}
\begin{tabular}{@{}lc@{}}
\toprule
Split Strategy & $E_{\text{split}}$ \\
\midrule
Low-$z$ train, High-$z$ test ($z = 1$ threshold) & 1.4 \\
Alternating bins & 2.1 \\
Random 50/50 (averaged over 10 trials) & $1.2 \pm 0.8$ \\
\bottomrule
\end{tabular}
\end{table}

All data-split e-values yield $E < 3$, consistent with the power calibration showing median $E \approx 2.7$ under $H_1$ (Section~\ref{sec:datasplit}).

\paragraph{Uniform mixture e-values.} A uniform mixture e-value---which averages the likelihood ratio over a grid of $(w_0, w_a)$ values---yields $E \approx 15$ with the default prior range, and varies from approximately 15 to 97 as the prior range is narrowed. On conventional scales, $E = 15$ ($\ln E \approx 2.7$) constitutes moderate-to-strong evidence. The variation from 15 to 97 reflects the Occam razor operating as expected: a narrower prior that concentrates mass near the MLE produces larger e-values because it ``wastes'' less probability on parameter values far from the data. This is a standard feature of Bayesian and mixture-based tests, not a deficiency.

\paragraph{Convergence of valid methods.} Two independent valid e-value methods give consistent results:
\begin{equation}
\text{Uniform mixture:} \quad E \approx 15 \quad (\ln E \approx 2.7) \qquad \text{LOO average:} \quad E \approx 10 \quad (\ln E \approx 2.3)
\end{equation}
This convergence on $E \approx 10$--$15$ (moderate evidence) is reassuring, as the two methods are constructed differently: the mixture e-value averages over a prior grid using all data, while the LOO average uses out-of-sample prediction. The data-split $E = 1.4$ is consistent with these once its limited power is accounted for (median $E_{\text{split}} \approx 2.7$ under $H_1$).

\subsection{Power-Optimized Point-Alternative E-Value}
\label{sec:ump}

Following the construction described in Section~\ref{sec:ump_theory} and motivated by a suggestion of P.~Gr\"unwald (private communication), we compute the e-value associated with a single, pre-specified point alternative $\theta_1^*$ chosen to maximize power against a representative ensemble of $w_0w_a$CDM alternatives. We search over a $15 \times 15$ grid of candidates $(w_0, w_a) \in [-1.5, -0.5] \times [-2.0, 1.0]$, evaluating power via 500 Monte Carlo simulations under four representative ``true'' alternatives: the DESI BAO best-fit $(-0.75, -1.05)$, a milder evolution $(-0.85, -0.50)$, a DES-Y5-like model $(-0.65, -1.20)$, and a Pantheon+-like model $(-0.90, -0.20)$. We use $\alpha = 0.01$ (threshold $1/\alpha = 100$) and refine the grid optimum via Nelder--Mead optimization with 1000 simulations per evaluation.

The selected point alternative is $\theta_1^* = (w_0^* = -0.87,\, w_a^* = -0.45)$, with minimum power $72\%$ and average power $84\%$ across the four representative alternatives at $\alpha = 0.01$. Evaluating the simple likelihood ratio on the real DESI data:
\begin{equation}
E^{*} = \frac{L(\text{data} \mid w_0^*, w_a^*)}{L(\text{data} \mid \text{\LCDM})} \approx 2180 \quad (\Delta\chi^2 \approx 15.4, \;\; \ln E^{*} \approx 7.7).
\label{eq:ump}
\end{equation}

\paragraph{Validity.} Because $\theta_1^*$ is determined entirely by the simulation procedure---which depends on the covariance structure and on the pre-specified ensemble of true alternatives, but \emph{not} on the observed data---the resulting likelihood ratio is a valid e-value by Proposition~\ref{prop:lr}: $\E[E^{*} \mid H_0] = 1$ exactly. We verify this empirically: under 500 simulations with \LCDM{} as the true data-generating model, $\hat{\E}[E^{*}] = 0.62$ and $\hat{P}(E^{*} \geq 100) = 0$. The validity of $E^{*}$ does not depend on the choice of $\alpha$ or sample size; these enter only via the optimization step that selects $\theta_1^*$.

\paragraph{Scope of the result.} The numerical value $E^{*} \approx 2180$ is large and the e-value is valid, but its scope must be stated precisely. $E^{*}$ tests \LCDM{} against the \emph{single point} $\theta_1^*$, not against the composite $w_0w_a$CDM family. It tells us that the data are $\sim 2180$ times more probable under this one specific $(w_0, w_a)$ than under \LCDM, where $\theta_1^*$ was chosen in advance as a good ``bet'' against alternatives near the DESI best fit. It does \emph{not} answer the broader question---how strongly do the data favor the composite $w_0w_a$CDM family?---which is what the mixture and LOO e-values are designed to address. Those methods integrate over uncertainty about which $(w_0, w_a)$ to test and unavoidably pay an averaging penalty; the point-alternative construction avoids that penalty by committing to a single value.

The two answers are therefore complementary, not contradictory: the mixture and LOO methods report moderate evidence ($E \approx 10$--$15$) against the composite alternative, whereas $E^{*}$ reports strong evidence against the specific point $\theta_1^*$. Both are valid; reporting both makes transparent how much of the spread between the mixture/LOO numbers and the (invalid) maximized likelihood ratio is attributable to the averaging penalty rather than to weakness of the underlying signal.

\paragraph{Dependence on the simulation design.} The point $\theta_1^*$ depends on the ensemble of ``true'' alternatives used in the power simulation and on the choice of $\alpha$; a different ensemble would generally select a different $\theta_1^*$ and yield a different $E^{*}$ on the real data. This is analogous to the prior dependence of Bayesian model selection and of the uniform mixture e-value, and should be acknowledged whenever the numerical value is quoted: $E^{*} \approx 2180$ is the e-value attached to one particular pre-specified design, not a unique number associated with the data. In supplementary checks, reasonable variations in the simulation ensemble move $\theta_1^*$ over a region whose corresponding $E^{*}$ values lie in $\sim 10^{2}$--$10^{3}$ on the observed DESI data, so the qualitative conclusion is robust even though the precise number is not.

\paragraph{Position relative to the headline.} Despite its size, we do not regard $E^{*}$ as the headline result. The most informative novel finding remains the cross-dataset tension of Section~\ref{sec:cross}: DES-Y5's best-fit $(w_0, w_a)$ predict DESI \emph{worse} than \LCDM{} ($E = 0.19$) while Pantheon+'s predict it well ($E = 2049$). This is a statement about whether DESI and the supernova catalogs tell a coherent story about dark energy, and it cannot be made to go away by clever choice of alternative. $E^{*}$ contributes a different kind of insight---that, within the DESI BAO data alone, the e-value framework is not intrinsically incapable of registering strong evidence, and that the moderate mixture/LOO numbers reflect honest averaging rather than a verdict of weakness.

\subsection{Same-Data Likelihood Ratio (Invalid as E-Value)}
\label{sec:samedata}

For reference, fitting $(w_0, w_a)$ to all 13 DESI measurements and computing the maximized likelihood ratio on the same data yields:
\begin{equation}
\text{LR}_{\text{same}} \approx 3300 \quad (\Delta\chi^2 \approx 16.2)
\end{equation}

\textbf{This is not an e-value.} As shown in Section~\ref{sec:evalues}, plugging in the MLE gives $\E_{H_0}[\text{LR}] = \infty$ for $k=2$ extra parameters, violating the defining property $\E[E \mid H_0] \leq 1$. We include this maximized likelihood ratio solely to illustrate the contrast with the valid cross-dataset and data-split e-values presented above.

%=====================================================================
\section{Discussion}
\label{sec:discussion}
%=====================================================================

\subsection{Relation to Existing Literature}

Our findings complement two independent analyses:

\begin{enumerate}
    \item \textbf{Bayesian model comparison} \citep{OngBayesian2025}: Using nested sampling, these authors found that the Bayes factor for DESI+CMB modestly favors \LCDM{} ($\ln \mathcal{B} = -0.57 \pm 0.26$), and attributed the frequentist significance to a tension between DESI and DES-Y5 that $w_0w_a$CDM resolves. Our cross-dataset e-value ($E = 0.19$ for DES-Y5$\rightarrow$DESI) provides independent evidence of this tension. Following \citet{GrunwaldSafeTesting2024}, we note that because our null is simple, the Bayes factor of \citet{OngBayesian2025} is itself a valid e-value: the prior-mixture marginal likelihood divided by the simple-null likelihood is precisely an instance of the mixture e-value of Section~\ref{sec:evalues}, with the prior playing the role of the mixture distribution. The result $\ln \mathcal{B} = -0.57$ thus admits both a Bayesian and an e-value reading, and its prior dependence is the same averaging dependence we observe in our uniform mixture e-values (Section~\ref{sec:mixture}). The additional methods explored here are valuable not because they resolve this prior dependence but because they show how the answer moves as one varies the treatment of the composite alternative, from a fully averaged Bayes-factor-style mixture ($\ln \mathcal{B} = -0.57$) to a single pre-specified point ($E^{*} \approx 2180$).

    \item \textbf{Tension diagnostics} \citep{WangMota2025}: These authors documented parameter inconsistencies between CMB, DESI, and supernova catalogs. Our observation that different supernova catalogs yield very different predictive performance on DESI (Pantheon+: $E = 2049$ vs DES-Y5: $E = 0.19$) aligns with their conclusion that combining these datasets is problematic.
\end{enumerate}

All three approaches---Bayesian model comparison, tension metrics, and e-value analysis---converge on the conclusion that inter-dataset tensions are a central issue. Within the DESI data alone, our two valid e-value methods converge on moderate evidence ($E \approx 10$--$15$), while the Bayes factor suggests weak evidence for \LCDM{} and the data-split e-value is inconclusive due to limited power. This spread is informative: it indicates that the answer depends on how one handles the composite alternative hypothesis, though the convergence of the mixture and LOO methods narrows the range considerably.

\subsection{Comparison with Information Criteria}
\label{sec:infocriteria}

To place the e-value results in the context of the full spectrum of evidence measures, we report the standard information criteria for the \LCDM{} vs.\ $w_0w_a$CDM comparison. The $w_0w_a$CDM model has $k = 2$ extra parameters relative to \LCDM, and the $\chi^2$ improvement is $\Delta\chi^2 = 11.9$.

\paragraph{AIC.} The Akaike Information Criterion \citep{BurnhamAnderson2002} penalizes by $2k$:
\begin{equation}
\Delta\text{AIC} = \Delta\chi^2 - 2k = 11.9 - 4 = 7.9
\end{equation}
On the \citet{BurnhamAnderson2002} scale, $\Delta\text{AIC} > 6$ constitutes ``strong'' evidence favoring the more complex model. Thus, AIC strongly favors $w_0w_a$CDM.

\paragraph{BIC.} The Bayesian Information Criterion \citep{Schwarz1978} penalizes by $k \ln n$:
\begin{equation}
\Delta\text{BIC} = \Delta\chi^2 - k \ln n = 11.9 - 2 \ln 13 \approx 11.9 - 5.13 = 6.8
\end{equation}
On the \citet{KassRaftery1995} scale, $6 < \Delta\text{BIC} < 10$ constitutes ``strong'' evidence. Thus, BIC also favors $w_0w_a$CDM.

\paragraph{Caveat on the effective sample size $n$.} We used $n = 13$ (the number of BAO summary statistics) in the BIC formula, but the meaning of $n$ in cosmological model selection is ambiguous \citep{Liddle2007}. If one instead treats the underlying galaxy pairs (millions) as the effective sample, $\ln n \sim 30$ and $\Delta\text{BIC}$ becomes strongly \textit{negative}, favouring \LCDM. This is a well-known limitation of BIC in cosmological applications.

\paragraph{Why different measures disagree.} The evidence landscape for \LCDM{} vs.\ $w_0w_a$CDM depends on how one penalizes model complexity:

\begin{table}[H]
\centering
\caption{Summary of evidence measures for $w_0w_a$CDM vs.\ \LCDM}
\label{tab:evidence_summary}
\begin{tabular}{@{}llll@{}}
\toprule
\textbf{Method} & \textbf{Value} & \textbf{Favors} & \textbf{Complexity penalty} \\
\midrule
Frequentist $p$-value ($k=2$) & $p = 0.0026$ ($2.8\sigma$) & $w_0w_a$CDM & None (likelihood only) \\
$\Delta$AIC & 7.9 & $w_0w_a$CDM & $2k = 4$ \\
$\Delta$BIC ($n=13$) & 6.8 & $w_0w_a$CDM & $k\ln n \approx 5.1$ \\
Bayes factor (Ong et al.) & $\ln\mathcal{B} = -0.57$ & \LCDM & Full prior volume \\
E-value (mixture, default) & $E \approx 15$ & $w_0w_a$CDM & Prior grid averaging \\
E-value (LOO average) & $E \approx 10$ & $w_0w_a$CDM & Out-of-sample (avg) \\
E-value (point alt., power-opt.) & $E^{*} \approx 2180$ & single $\theta_1^*$ vs.\ \LCDM & None (commits to a point) \\
E-value (data-split) & $E = 1.4$ & Inconclusive & Out-of-sample (underpowered) \\
\bottomrule
\end{tabular}
\end{table}

\noindent AIC includes no Occam factor (it penalizes only the number of parameters), so it tends to favor more complex models. BIC approximately includes an Occam factor through $k\ln n$, but the result depends on what counts as a ``data point.'' The Bayes factor includes the full Occam penalty through the prior volume: with broad priors over $(w_0, w_a)$, much of the prior mass falls on parameter values that fit the data poorly, penalizing $w_0w_a$CDM. The mixture and LOO e-values ($E \approx 10$--$15$) are intermediate between the information criteria and the Bayes factor: like the Bayes factor, they pay an averaging penalty over the alternative parameter space; unlike the Bayes factor, they do not require a normalized prior and admit a frequentist guarantee for any sample size. The point-alternative e-value $E^{*} \approx 2180$ has the smallest implicit ``averaging penalty'' of any of the valid methods because it commits in advance to a single $\theta_1^*$, but its scope is correspondingly narrower: it tests \LCDM{} against this one point, not against the composite model as a whole (Section~\ref{sec:ump}). The data-split e-value ($E = 1.4$) is inconclusive due to limited power rather than reflecting weak evidence.

These different answers are not contradictory---they reflect different questions. AIC asks which model will predict future data better on average (favoring $w_0w_a$CDM). BIC approximates the posterior model probability (ambiguous). The Bayes factor asks how much the data update the relative prior odds between two models (favoring \LCDM{} with broad priors over $(w_0, w_a)$). The mixture and LOO e-values ask, ``averaging honestly over our uncertainty about which $(w_0, w_a)$ to test, how much do the data favor the alternative?'' The point-alternative e-value asks the narrower but more powerful question, ``how much do the data favor a single specific $(w_0, w_a)$ chosen up-front for power?'' All are valid; they answer related but different questions, and their numerical spread (from $E \approx 10$ to $E^{*} \approx 2180$) is informative about the cost of how the composite alternative is summarized.

\subsection{Frequentist Significance for $k = 2$ Parameters}
\label{sec:sigma}

For readers unfamiliar with the conversion, we briefly review the proper frequentist significance for $\Delta\chi^2$ with $k = 2$ extra parameters. Under the null hypothesis (\LCDM), $\Delta\chi^2$ follows a $\chi^2$ distribution with 2 degrees of freedom (by Wilks' theorem). The $p$-value is:
\begin{equation}
p = 1 - F_{\chi^2}(11.9;\, 2) = e^{-11.9/2} = e^{-5.95} \approx 0.0026
\end{equation}
where we used the exact CDF for the $\chi^2(2)$ distribution, $F_{\chi^2}(x;\, 2) = 1 - e^{-x/2}$. Converting to a two-sided Gaussian equivalent:
\begin{equation}
\sigma = \Phi^{-1}(1 - p/2) \approx \Phi^{-1}(0.9987) \approx 2.8\sigma
\end{equation}

DESI's own analysis correctly performs this conversion using the $\chi^2(2)$ distribution and Wilks' theorem (their equation~22 in \citealt{DESI2025}), reporting BAO-only significance of approximately $3.0\sigma$ (for their slightly different $\Delta\chi^2$). The naive formula $\sigma = \sqrt{\Delta\chi^2}$, which is valid only for $k = 1$, would overestimate the significance; we mention it only because it sometimes appears in informal discussions of the result.

\paragraph{Converting e-values to a $\sigma$-scale.} The mapping $\sigma \approx \sqrt{2\ln E}$ that we use in places to compare e-values to frequentist significances is approximate: it is exact for a one-parameter likelihood ratio under Gaussian Wilks asymptotics, and it is the natural scale on which the simple-vs-simple e-value $E^{*}$ can be compared to a point-null $z$-test, but $e$-values and $p$-values answer different questions and there is no exact universal conversion between them. The mixture and LOO averages, in particular, should be read on this scale only as a rough comparison aid. Concretely, our averaging methods correspond to $E \approx 10$--$15$ ($\approx 2.2$--$2.4\sigma$) versus DESI's BAO-only $\approx 3.0\sigma$: roughly $0.7\sigma$ of the frequentist signal is absorbed by the Occam penalty over $(w_0, w_a)$ when one is honest about uncertainty in the alternative.

\subsection{$w_0w_a$CDM as a Tension Absorber}

The cross-dataset analysis (Section~\ref{sec:cross}) raises a deeper question about the nature of the reported 3--4$\sigma$ evidence. The $w_0w_a$CDM model introduces two free parameters ($w_0$ and $w_a$) that can adjust distance-redshift relations to better fit the data. When multiple datasets pull in different directions---as the ${\sim}10{,}000\times$ Pantheon+/DES-Y5 asymmetry demonstrates---these extra parameters can absorb the tension by finding a compromise fit that improves $\chi^2$ overall without corresponding to any physical dark energy dynamics. In this scenario, $w_0w_a$CDM acts as a \textit{tension absorber} rather than a detector of new physics.

This interpretation is supported by the pattern observed in the cross-prediction matrix: the DESI+CMB+SNe combination yields 3--4$\sigma$ preference for $w_0w_a$CDM, yet the constituent datasets disagree on \textit{which} $(w_0, w_a)$ values are preferred. DES-Y5 favours $(w_0, w_a) = (-0.65, -1.20)$, incompatible with DESI ($E = 0.19$), while Pantheon+'s $(-0.90, -0.20)$ is compatible ($E = 2049$). The combined ``significance'' thus reflects the model's ability to split the difference between discrepant datasets rather than a coherent signal on which all experiments converge---the same mechanism identified by \citet{OngBayesian2025}, who showed that the 2.95$\sigma$ DESI/DES-Y5 tension within \LCDM{} is precisely what $w_0w_a$CDM resolves. Until the source of the Pantheon+/DES-Y5 discrepancy is understood, the 3--4$\sigma$ frequentist significance may be measuring the magnitude of dataset inconsistency rather than evidence for new physics.

\subsection{What E-Values Add}

The e-value perspective contributes:
\begin{enumerate}
    \item \textbf{Direct test of generalization:} Data-splitting and LOO cross-validation directly ask whether fitted parameters predict held-out data, addressing overfitting concerns.

    \item \textbf{Multiple valid methods spanning a known range:} The uniform mixture ($E \approx 15$), LOO average ($E \approx 10$), and point-alternative e-value $E^{*} \approx 2180$ are constructed differently and answer subtly different questions, but all satisfy $\E[E \mid H_0] \leq 1$. The mixture and LOO numbers represent the answer when one averages honestly over uncertainty in the alternative; $E^{*}$ represents the answer when one commits in advance to a single alternative chosen for power. Reporting both ends of the spectrum makes transparent how much of the spread is attributable to the averaging penalty rather than to the underlying signal being weak.

    \item \textbf{Different form of prior dependence:} Data-split and LOO e-values do not require an explicit prior distribution over $(w_0, w_a)$. However, they are not prior-free: using the point MLE from the training set is effectively equivalent to a point-mass prior centered at $\hat{\theta}$, and results depend on the split or fold choice. The uniform mixture e-value explicitly integrates over a prior grid and is transparent about this dependence.

    \item \textbf{Interpretable metric:} An e-value of $E$ directly means ``the data are $E$ times more consistent with the alternative than expected under the null.''
\end{enumerate}

\subsection{Limitations of This Analysis}
\label{sec:limitations}

\begin{enumerate}
    \item \textbf{Data-split power and split dependence.} The redshift-split test is underpowered for $w_a$ for the reasons explained in Section~\ref{sec:datasplit}; results also vary modestly across alternative splits (Table~\ref{tab:splits}), though all yield $E < 3$. We rely on the LOO average and uniform mixture, which use all the data, for our within-dataset inferences.

    \item \textbf{Simplified cosmology and no nuisance-parameter marginalization.} We fix $\{h, \Omega_m, r_d\}$ to Planck 2018 values and use closed-form distance calculations rather than full Boltzmann codes (CAMB/CLASS). A full treatment would marginalize over these parameters, with two competing effects on the e-values: broadening the effective alternative reduces the uniform mixture (more ``wasted'' prior volume), while allowing the fit to adjust $\Omega_m$ and $h$ may improve $\chi^2$ and increase it. This limitation is shared with any analysis that conditions on external parameter constraints, including standard $\Delta\chi^2$ and BIC.

    \item \textbf{Point estimates for the supernova constraints.} We use published best-fit $(w_0, w_a)$ values rather than full posteriors. A more rigorous treatment would marginalize the cross-dataset e-value over the supernova posterior, accounting for parameter uncertainty in the training dataset.
\end{enumerate}

%=====================================================================
\section{Conclusion}
\label{sec:conclusion}
%=====================================================================

We applied e-value methodology to the DESI DR2 BAO data to examine the evidence for $w_0w_a$CDM from several angles. The picture that emerges is methodology-sensitive rather than unambiguous in either direction.

\begin{enumerate}
    \item \textbf{Cross-dataset tension is the headline finding.} DES-Y5's best-fit $w_0w_a$CDM parameters predict the DESI BAO data \emph{worse} than \LCDM{} ($E = 0.19$), while Pantheon+'s parameters predict it well ($E = 2049$). The ${\sim}10{,}000\times$ asymmetry between the two leading supernova catalogs is a direct, calibrated detection of inter-dataset inconsistency: if $w_0w_a$CDM described real physics, the catalogs should agree on which $(w_0, w_a)$ region is preferred, and they do not. This raises the possibility that $w_0w_a$CDM is absorbing inter-dataset tension rather than detecting new physics, and aligns with the parallel findings of \citet{OngBayesian2025} and \citet{WangMota2025}.

    \item \textbf{Information criteria favor $w_0w_a$CDM, the Bayes factor favors \LCDM, and valid e-value methods span the gap.} $\Delta\text{AIC} = 7.9$ and $\Delta\text{BIC} = 6.8$ constitute ``strong'' evidence for $w_0w_a$CDM on standard scales, while the Ong et al.\ Bayes factor modestly favors \LCDM. The divergence reflects different penalties for the alternative parameter space; the BIC value is also subject to the well-known ambiguity in the effective sample size for cosmological data \citep{Liddle2007}.

    \item \textbf{Within the DESI BAO data alone, valid e-value methods span moderate to strong evidence depending on how the composite alternative is summarized.} The uniform mixture ($E \approx 15$) and LOO average ($\bar{E} \approx 10$) average over uncertainty about which $(w_0, w_a)$ to test and yield moderate evidence ($\approx 2.2$--$2.4\sigma$). The power-optimized point-alternative e-value $E^{*} \approx 2180$, computed via a construction suggested by P.~Gr\"unwald, commits in advance to a single $(w_0, w_a)$ chosen for power against alternatives near the DESI best fit, and so avoids the averaging penalty entirely. These numbers answer related but different questions and are complementary, not competing: the spread between them quantifies the cost of being explicit about uncertainty in the alternative. The redshift-split e-value $E = 1.4$ is, by power calibration, inconclusive rather than negative and does not weigh against the other valid methods.
\end{enumerate}

The most informative novel finding remains the cross-dataset tension: until the source of the Pantheon+/DES-Y5 discrepancy is understood, the 3--4$\sigma$ frequentist signal from the combined dataset should be interpreted with caution, and within the DESI data alone the evidence for dynamical dark energy is moderate, not yet at the level of a discovery. Future data from DESI DR3+, Euclid, and the Roman Space Telescope will be essential for separating new physics from systematic effects.

%=====================================================================
\section*{Data and Code Availability}
%=====================================================================

Analysis code and data files are available at:
\begin{center}
\url{https://github.com/jinyoungkim927/desi-evalue-paper}
\end{center}

DESI BAO data from: \url{https://github.com/CobayaSampler/bao_data}

%=====================================================================
\section*{Acknowledgments}
%=====================================================================

We thank Peter Gr\"unwald (CWI Amsterdam) for detailed feedback on the e-value methodology and for suggesting the power-optimized point-alternative construction discussed in Section~\ref{sec:ump}, and Ruodu Wang (University of Waterloo) for confirming the validity of the leave-one-out averaging argument and the convergence interpretation. DFM and AT thank the Research Council of Norway for their support and the resources provided by UNINETT Sigma2---the National Infrastructure for High-Performance Computing and Data Storage in Norway. We also thank the DESI collaboration for the public release of the DR2 BAO data products used in this analysis.

%=====================================================================
\appendix
\section{E-Value Theory}
\label{app:proofs}
%=====================================================================

This appendix provides formal statements of the e-value properties used in the main text.

\begin{definition}[E-Value]
A random variable $E \geq 0$ is an \textbf{e-value} for testing $H_0$ if $\E[E \mid H_0] \leq 1$.
\end{definition}

\begin{theorem}[Ville's Inequality]
\label{thm:ville}
If $E$ is an e-value, then for any $\alpha \in (0,1)$:
\begin{equation}
\Prob_{H_0}(E \geq 1/\alpha) \leq \alpha
\end{equation}
\end{theorem}

\begin{proof}
By Markov's inequality: $\Prob(E \geq 1/\alpha) \leq \alpha \cdot \E[E] \leq \alpha$.
\end{proof}

\begin{proposition}[Likelihood Ratio E-Value]
\label{prop:lr}
For simple hypotheses $H_0: P = P_0$ versus $H_1: P = P_1$, the likelihood ratio $E = P_1(X)/P_0(X)$ satisfies $\E_{P_0}[E] = 1$.
\end{proposition}

\begin{proof}
$\E_{P_0}[E] = \int \frac{P_1(x)}{P_0(x)} P_0(x) \, dx = \int P_1(x) \, dx = 1$.
\end{proof}

\begin{proposition}[Data-Split E-Value Validity]
\label{prop:split}
Let $D = (D_{\text{train}}, D_{\text{test}})$ be independent data splits. Let $\hat{\theta} = \hat{\theta}(D_{\text{train}})$ be parameters fitted to the training data. Then:
\begin{equation}
E = \frac{L(D_{\text{test}} \mid \hat{\theta})}{L(D_{\text{test}} \mid H_0)}
\end{equation}
is a valid e-value for testing $H_0$.
\end{proposition}

\begin{proof}
Conditional on $D_{\text{train}}$, $\hat{\theta}$ is fixed. Under $H_0$, $D_{\text{test}} \sim P_0$ independent of $D_{\text{train}}$. Thus:
\begin{equation}
\E[E \mid D_{\text{train}}] = \int \frac{L(x \mid \hat{\theta})}{L(x \mid H_0)} P_0(x) \, dx = 1
\end{equation}
by the same argument as Proposition~\ref{prop:lr}. Taking expectations: $\E[E] = \E[\E[E \mid D_{\text{train}}]] = 1$.
\end{proof}

\begin{proposition}[Uniform Mixture E-Value Validity]
\label{prop:mixture}
Let $\mathcal{G} = \{\theta_1, \ldots, \theta_M\}$ be a finite set of parameter values specified before seeing the data. The uniform mixture
\begin{equation}
E_{\text{mix}} = \frac{1}{M} \sum_{j=1}^M \frac{L(\text{data} \mid \theta_j)}{L(\text{data} \mid H_0)}
\end{equation}
is a valid e-value.
\end{proposition}

\begin{proof}
Each term $E_j = L(\text{data} \mid \theta_j) / L(\text{data} \mid H_0)$ is a likelihood ratio for a simple alternative, so $\E_{H_0}[E_j] = 1$ by Proposition~\ref{prop:lr}. By linearity of expectation:
\begin{equation}
\E_{H_0}[E_{\text{mix}}] = \frac{1}{M} \sum_{j=1}^M \E_{H_0}[E_j] = \frac{1}{M} \cdot M = 1.
\end{equation}
\end{proof}

\begin{proposition}[LOO Average E-Value Validity]
\label{prop:loo}
Let $E_1, \ldots, E_K$ be leave-one-out e-values, where $E_k = L(D_k \mid \hat{\theta}_{-k}) / L(D_k \mid H_0)$ and $\hat{\theta}_{-k}$ is fitted on all data except fold $k$. The average $\bar{E} = \frac{1}{K}\sum_{k=1}^K E_k$ is a valid e-value.
\end{proposition}

\begin{proof}
Each $E_k$ is a valid e-value by Proposition~\ref{prop:split}: conditional on the training data $D_{-k}$, $\hat{\theta}_{-k}$ is fixed, and the held-out data $D_k$ provides an independent test. Therefore $\E[E_k \mid H_0] \leq 1$ for each $k$. By linearity:
\begin{equation}
\E[\bar{E} \mid H_0] = \frac{1}{K}\sum_{k=1}^K \E[E_k \mid H_0] \leq \frac{1}{K} \cdot K = 1.
\end{equation}
Note: the \textit{product} $\prod_k E_k$ does \textit{not} have this guarantee because the $E_k$ are not independent (the training sets overlap).
\end{proof}

\begin{remark}[Gaussian Likelihoods]
For multivariate Gaussian likelihoods with known covariance $C$:
\begin{equation}
E = \exp\left(\frac{\chi^2(H_0) - \chi^2(\hat{\theta})}{2}\right)
\end{equation}
where $\chi^2(\theta) = (d - t(\theta))^T C^{-1} (d - t(\theta))$.
\end{remark}

%=====================================================================
\bibliographystyle{plainnat}
\begin{thebibliography}{99}

\bibitem[Burnham \& Anderson(2002)]{BurnhamAnderson2002}
Burnham, K. P., \& Anderson, D. R. 2002, \textit{Model Selection and Multimodel Inference: A Practical Information-Theoretic Approach}, 2nd ed. (Springer)

\bibitem[Chevallier \& Polarski(2001)]{Chevallier2001}
Chevallier, M., \& Polarski, D. 2001, ``Accelerating Universes with Scaling Dark Matter,'' International Journal of Modern Physics D, 10, 213

\bibitem[DESI Collaboration(2025)]{DESI2025}
DESI Collaboration 2025, ``DESI DR2 Results II: Measurements of Baryon Acoustic Oscillations and Cosmological Constraints,'' arXiv:2503.14738

\bibitem[Gr{\"u}nwald(2024)]{GrunwaldBeyondNeymanPearson}
Gr{\"u}nwald, P. 2024, ``Beyond Neyman--Pearson: E-Values Enable Hypothesis Testing with a Data-Driven Alpha,'' Proceedings of the National Academy of Sciences, 121, e2302098121

\bibitem[Kass \& Raftery(1995)]{KassRaftery1995}
Kass, R. E., \& Raftery, A. E. 1995, ``Bayes Factors,'' Journal of the American Statistical Association, 90, 773

\bibitem[Gr{\"u}nwald et al.(2024)]{GrunwaldSafeTesting2024}
Gr{\"u}nwald, P., de Heide, R., \& Koolen, W. M. 2024, ``Safe Testing,'' Journal of the Royal Statistical Society Series B, 86, 1091--1128

\bibitem[Liddle(2007)]{Liddle2007}
Liddle, A. R. 2007, ``Information Criteria for Astrophysical Model Selection,'' MNRAS, 377, L74

\bibitem[Linder(2003)]{Linder2003}
Linder, E. V. 2003, ``Exploring the Expansion History of the Universe,'' Physical Review Letters, 90, 091301

\bibitem[Ong et al.(2025)]{OngBayesian2025}
Ong, D. D. Y., Yallup, D., \& Handley, W. 2025, ``A Bayesian Perspective on Evidence for Evolving Dark Energy,'' arXiv:2511.10631

\bibitem[Ramdas et al.(2023)]{Ramdas2023}
Ramdas, A., Gr{\"u}nwald, P., Vovk, V., \& Shafer, G. 2023, ``Game-Theoretic Statistics and Safe Anytime-Valid Inference,'' Statistical Science, 38, 576

\bibitem[Schwarz(1978)]{Schwarz1978}
Schwarz, G. 1978, ``Estimating the Dimension of a Model,'' Annals of Statistics, 6, 461

\bibitem[Shafer(2021)]{Shafer2021}
Shafer, G. 2021, ``Testing by Betting: A Strategy for Statistical and Scientific Communication,'' Journal of the Royal Statistical Society Series A, 184, 407

\bibitem[Vovk \& Wang(2021)]{Vovk2021}
Vovk, V., \& Wang, R. 2021, ``E-values: Calibration, Combination, and Applications,'' Annals of Statistics, 49, 1736

\bibitem[Wang \& Mota(2025)]{WangMota2025}
Wang, D., \& Mota, D. 2025, ``Did DESI DR2 Truly Reveal Dynamical Dark Energy?,'' arXiv:2504.15222

\end{thebibliography}

\end{document}
